\documentclass[11pt]{article}

% The debug report: a first person postmortem grouped by theme, built from the
% engineering log. It is an equal rank deliverable, not an afterthought, and I do
% not pad it. Entries appear here once the problem behind them is understood and
% fixed, so this document trails the engineering log by exactly the lag between
% hitting a problem and resolving it.

\usepackage[margin=1in]{geometry}
\usepackage{amsmath}
\usepackage{booktabs}
\usepackage{xcolor}
\usepackage[hidelinks]{hyperref}

\title{GPU Roofline Profiler: Debug Report}
\author{Olajide Badejo}
\date{\today}

\begin{document}

% Same hand built title page as the main report: title centred both ways, byline
% at the foot. \end{titlepage} already starts a new page, so no \clearpage here.
\begin{titlepage}
  \centering
  \vspace*{\fill}
  {\LARGE\bfseries GPU Roofline Profiler: Debug Report\par}
  \vspace*{\fill}
  {\large Olajide Badejo\par}
  \vspace{0.9em}
  {\large \today\par}
  \vspace{1.5cm}
\end{titlepage}

\begin{abstract}
This is the honest account of what went wrong while building the profiler and
how I fixed it, grouped by theme rather than by date. The dated raw material
lives in \texttt{docs/\allowbreak ENGINEERING\_LOG.md}; this document is the
postmortem that draws the lessons out of it.

A theme runs through most of these. The failures that cost the most time were not
the ones that produced an error message. They were the ones that produced a
plausible looking number: a profiler that collected nothing and said nothing, a
manifest that no parser could read, a metric that flattered the wrong thing, an
optimisation that quietly targeted the wrong GPU architecture.
\end{abstract}
\clearpage

\tableofcontents
\clearpage

\section{Environment and tooling}

\subsection{The machine had the GPU but none of the toolchain}

The first probe of the target machine was a surprise in the honest direction.
\texttt{nvidia-smi} reported a healthy RTX 5070 on driver 610.62 with a CUDA user
mode driver at version 13.3, so the card and its runtime were fine. Everything
needed to actually build against it was missing: no CUDA Toolkit and therefore no
\texttt{nvcc}, no Visual Studio and therefore no \texttt{cl.exe} for \texttt{nvcc}
to use as its host compiler, and neither Nsight Systems nor Nsight Compute. The
machine had the gaming and display stack but not the developer stack.

The root cause was mundane: the tools were simply never installed. The right
response was to resolve the toolchain deliberately rather than let large
installers needing elevated permissions run unattended. In the meantime
I built every part of the project that does not need a compiler, so the wait was
not idle: the repository scaffold, the dash lint, the Python analysis package
with its tests, and the report skeletons all came first.

\subsection{The Python pins predated the interpreter}

The machine's usable Python is CPython 3.14. The dependency versions I first
pinned were older, and pip had no 3.14 wheels for them, so it fell back to
building pandas from source with meson, which then failed for the same reason
\texttt{nvcc} would have: no MSVC. The fix was to install the current releases
that do ship 3.14 wheels and to re-pin \texttt{requirements.txt} to exactly those
versions. Pinning to what actually installs and passes on this box is more honest
than pinning to an aspiration that cannot be built here.

\subsection{CUDA 13 removed the clock properties}

A trivial smoke test would not compile: \texttt{cudaDeviceProp} has no member
\texttt{clockRate}, and none named \texttt{memoryClockRate} either. CUDA 13
removed both. They were deprecated through 12.x and are now gone.

This mattered well beyond a smoke test, because the theoretical peak derivation
needs the SM clock and the memory clock, so the entire ceiling calculation rested
on two fields that no longer exist. The fix is to read them through
\texttt{cudaDeviceGetAttribute}, which still supports them, and to check the
returned status so that a future removal fails loudly rather than yielding a
silent zero and a peak of zero.

The check that the replacement is right: 48 SMs at 128 lanes is 6144 CUDA cores,
matching the vendor reference exactly, and at 2.625 GHz that gives 32.3 TFLOP/s
against a vendor quote near 31. The 192 bit bus at 14.001 GHz doubled for the
data rate gives 672.0 GB/s against a vendor quote of 672. Both landing on the
reference numbers is what says the formulas are right.

\subsection{The build silently targeted the wrong GPU}

The first successful CMake configure printed \texttt{CUDA architectures: 75}. The
card is sm\_120.

I had set \texttt{CMAKE\_CUDA\_ARCHITECTURES} inside an
\texttt{if(NOT DEFINED ...)} guard placed after \texttt{project()}. Enabling the
CUDA language in \texttt{project()} installs a default of its own, so by the time
my guard ran the variable was already defined and my detection never executed.

This is the quiet, expensive kind of bug. Nothing fails. The build just compiles
for an older architecture than the card, and every performance number that
followed would have been measured on code generated for the wrong target. The fix
is to choose the architecture before \texttt{project()}. I verified it afterwards
with \texttt{cuobjdump}, which now reports \texttt{sm\_120} SASS in every kernel
object rather than trusting the configure message a second time.

\subsection{A dependency that CMake 4 refuses to configure}

\texttt{FetchContent} pulled yaml-cpp 0.8.0 and configure failed: compatibility
with CMake below 3.5 has been removed, and this machine has CMake 4.3.1. I scoped
a \texttt{CMAKE\_POLICY\_VERSION\_MINIMUM} shim around that one dependency and
unset it immediately after, so the compatibility relaxation applies to the fetched
subproject and not to mine. Vendoring a patched copy would have been a
maintenance burden for the same result, and hand writing a YAML parser to dodge
one upstream version declaration would have been a poor trade.

\subsection{NVTX broke the build through windows.h}

Enabling NVTX turned a clean build into a syntax error in
\texttt{timing.hpp}, a file that had not changed and that has nothing to do with
profiling. \texttt{nvtx3/nvToolsExt.h} includes \texttt{windows.h}, which defines
\texttt{min} and \texttt{max} as preprocessor macros, and the offending line is
\texttt{std::max(1, std::min(wanted, max\_launches))}. Once those are macros the
expression stops parsing, and the error points at the victim rather than the
culprit. Defining \texttt{NOMINMAX} next to the include that causes the problem
fixes it and keeps the fix where the next person will look.

Alongside it, the CMake probe for NVTX tested \texttt{EXISTS} against
\texttt{CUDAToolkit\_INCLUDE\_DIRS}, which is a \emph{list} of two directories
rather than one path, so it silently reported "NVTX not found" on a machine where
the header was plainly present. \texttt{find\_path} searches the list properly.

\section{Measurement methodology}

\subsection{NVML explained a ceiling that looked impossible}

The measured FP32 peak came out at 33.5 TFLOP/s against a theoretical 32.26, or
104 percent of a number nothing is supposed to exceed.

The NVML trace settles it. The log records an SM clock averaging 2857 MHz and
peaking at 2865 MHz under load, while \texttt{cudaDevAttrClockRate} reports 2625
MHz. The runtime reports a nominal boost clock; the card boosts past it.
Recomputing at the observed clock gives 35.2 TFLOP/s, and the measured figure is
then 95 percent of it, exactly where a good FMA microbenchmark should land.

I changed nothing in the derivation. The theoretical ceiling stays derived from
the clock the runtime reports, because that is the reproducible definition, and
the report says plainly that the measured ceiling exceeds it and why. This is
precisely the case the NVML monitor exists for: without the clock trace this
would have looked like a counting bug and I would have gone hunting a factor of
two that was never there.

\subsection{Small kernels were measuring cache, not DRAM}

SAXPY at 4 million elements reported 1583 GB/s of achieved bandwidth against a
theoretical DRAM bandwidth of 672. At that size the two vectors total 32 MB and
fit inside this card's 48 MB L2, so after warmup the kernel never touches DRAM and
the figure is L2 bandwidth wearing a DRAM label. Past the cache the numbers
behave: 560 to 566 GB/s at 16 and 64 million elements.

The lesson is about the byte count rather than the timing. An achieved bandwidth
computed from a \emph{theoretical} byte count assumes every byte comes from DRAM,
and silently overstates DRAM traffic when it does not. This is the gap the whole
project is about, and it is why every point on the roofline is labelled with which
byte count produced its intensity.

\subsection{The shared memory tiling that did not pay}

I expected the naive GEMM to be memory bound and tiling to fix it. Across every
size measured, the tiled kernel is no faster than the naive one and at the
smaller sizes it is slower.

The counters explain it. Naive and tiled report L2 hit rates equal to within
measurement noise, and both move DRAM traffic at 5 to 8 GB/s against a 604 GB/s
ceiling, so neither is close to memory bound. The naive kernel would move 68.7 GB
if every operand read reached memory; it actually moves 80.8 MB, so cache absorbs
better than 99.8 percent of its redundant loads. Tiling exists to remove redundant
DRAM traffic and there is essentially none left to remove, while it still costs a
synchronisation per tile step. The optimisation targets a bottleneck this hardware
no longer has at this problem size.

What binds the naive kernel instead is instruction throughput: 98.5 percent of
peak sustained SM throughput while achieving under 2 TFLOP/s of useful arithmetic.
The fix is not to move data less far but to ask for it fewer times, which is what
register blocking does, and that rung is nearly four times faster.

I record this as the most useful result in the project. It is the one place where
measuring contradicted the textbook expectation, and the counters rather than the
narrative decided it.

\section{Kernel correctness}

\subsection{The tolerance was wrong, not the GEMM}

Every GEMM variant failed its correctness test at every size: naive, tiled,
register blocked, vectorized, and cuBLAS.

The clue that mattered was that all five reported the \emph{identical} error to
sixteen digits, and one of the five was cuBLAS. Five independent implementations
do not share a bug. They do share a reference and an error metric.

My metric was a per element relative error, dividing each deviation by
\texttt{max(|expected|, 1e-3)}. With random operands in $[-1, 1]$ each GEMM output
is a sum of $k$ signed products, so a good fraction of the entries land near zero
through cancellation, and dividing an ordinary absolute error by one of those
produces a huge ratio that measures the cancellation rather than the kernel.

Before changing anything I checked the arithmetic: a reported 3.2e-3 against the
1e-3 floor means an absolute error near 3.2e-6, while a 512 term sum of products
of values in $[-1, 1]$ has magnitude around 13. That is a true relative error near
2.5e-7, which is fp32 epsilon. The kernels were right the whole time.

The fix was a normwise relative error, the largest absolute deviation over the
largest magnitude in the reference, which is the standard way to check a GEMM and
does not lose meaning near zero. I deliberately did not simply loosen the
tolerance: that would have hidden the real problem behind a number chosen to make
tests pass, and would then have been too loose to catch an actual bug at large
$k$.

\subsection{The counters found a bug in my own kernels}

With the DRAM metrics finally collecting, the shared memory bank conflict counter
read 134.5 million for the register blocked GEMM and 93.4 million for the
vectorized one, against 40 thousand for the tiled transpose.

I had applied the classic \texttt{[TILE][TILE + 1]} padding to the transpose tile,
written a comment about why it matters, and then declared every GEMM shared tile
unpadded. The register blocked inner loop walks down a column of its tile, so the
address advances by a whole row per step; with an unpadded row length of 16 that
stride is a multiple of the 32 bank width and the warp serialises on one bank. I
had written the fix once and failed to apply it where it mattered most.

The fix carries a trap. The vectorized kernel reads its tile with a
\texttt{float4}, so row starts must stay 16 byte aligned; padding by one float
makes the row stride 260 bytes and turns an aligned vector load into undefined
behaviour rather than a slow path. Padding by four floats preserves alignment and
still moves each row off the colliding bank.

The measured outcome was not the clean sweep I expected. All 39 correctness tests
still pass, and at 4096 the vectorized kernel gained 18 percent, from 8746 to
10325 GFLOP/s, while the register blocked and tiled kernels did not move outside
run to run noise. Bank conflicts were the binding constraint in only one of the
three. I am reporting that rather than implying the optimisation helped
everywhere.

\section{Profiling tooling}

\subsection{Nsight Compute collected nothing, convincingly}

The first \texttt{ncu} run looked like a success. It attached to the process,
reported the kernel by name with its grid and block dimensions, ran the kernel,
and disconnected cleanly. Every metric value was \texttt{n/a}.

This is the most dangerous failure mode in the whole pipeline, because nothing
errors. A CSV full of \texttt{n/a} parses perfectly well, and piped straight into
the analysis it would have produced a report whose counter columns were empty for
a reason nobody had noticed.

The cause is that on Windows the NVIDIA driver restricts GPU performance counters
to administrators unless \texttt{RmProfilingAdminOnly} is set to 0 in the
registry. The value was not set at all, so the driver default applied, and the
shell was not elevated. This is the counter access caveat usually raised for
WSL2, and it bites on native Windows too.

The wrapper script now checks for elevation up front and refuses with an
explanation rather than producing a directory of \texttt{n/a}. I did not set the
registry value: it is a permanent, system wide relaxation of who may read GPU
counters, and that is the machine owner's decision to make deliberately rather
than a side effect of my wanting a profile. Nsight Systems, worth noting, needs no
elevation at all for CUDA and NVTX tracing.

\subsection{The metric names in the documentation do not exist here}

With elevation, most counters collected, and the two that mattered most did not:
\texttt{dram\_\_bytes\_\allowbreak read.sum} and
\texttt{dram\_\_bytes\_\allowbreak write.sum} still
returned \texttt{n/a}. Those are the names in most published examples. On this
chip the counters carry an \texttt{\_op\_} infix, and \texttt{ncu} does not reject
an unsupported metric name: it accepts it, profiles the kernel, and reports
\texttt{n/a}, producing a well formed CSV empty of the thing requested.

Reading \texttt{ncu -{}-query-metrics} on the machine, rather than trusting
documentation, gave the real names. The verified mapping is now recorded in
\texttt{docs/profiling\_guide.md} against the exact tool version, and the wrapper
counts \texttt{n/a} values after every pass and warns with the offending metric
names, so this specific silence can never happen twice.

\subsection{Two ways to be wrong about a percentage}

Folding the parsed metrics into per cell summaries, I aggregated everything the
same way. Counters and percentages are not the same kind of quantity: profiling
several launches and adding their byte counts gives the total traffic, which is
what we want, but adding two L2 hit rates of 97 percent does not give 194 percent
of anything. Extensive metrics now sum and intensive ones average.

That fix exposed a second thing, which is not a bug of mine. The naive GEMM
reports an L2 sector hit rate of 100.39 percent. A hit rate is hits over accesses
and cannot exceed 1. Re-running the same kernel gave 97.02, so it is noise in how
the hardware attributes sector hits.

I chose to keep the value verbatim and flag it rather than clamp it to 100.
Silently rewriting a measurement so it looks respectable is exactly the kind of
quiet dishonesty this project exists to avoid, and a reader who sees 100.39 with a
note attached learns something true about hardware counters that a tidy 100.00
would have hidden.

\section{Report pipeline}

\subsection{Underscores in a placeholder path broke the compile}

The main report guards every generated figure and table with a macro that prints a
visible placeholder when the artifact does not exist yet, so the document compiles
before and after the pipeline runs. The first version printed the missing path
inside \texttt{\textbackslash texttt}, and the paths contain underscores, which
LaTeX reads as math subscripts in text mode and refuses to typeset. Wrapping the
path in \texttt{\textbackslash detokenize} renders them literally.

\subsection{Every manifest was invalid JSON}

Generating the environment appendix from the run manifest threw a JSON decode
error: invalid escape. The driver wrote
\texttt{"config\_path": "configs\textbackslash sweep.yaml"} straight into the
document, and on Windows the path separator is a backslash, so
\texttt{\textbackslash s} is not a legal JSON escape and the whole file is
malformed. The device name went out unescaped too.

This mattered more than it looks. The project's own rule is that results without a
manifest do not exist, and an unparseable manifest is exactly as useless as a
missing one. Every run up to that point carried one. Nothing had noticed because
nothing had tried to read a manifest back until the report appendix needed it, and
the file looked perfectly reasonable to a human eye.

The fix is a \texttt{json\_escape} helper applied to every string written into the
manifest. The wider lesson: hand rolling JSON with a stream operator is fine for a
fixed schema right up until a value contains a character the format reserves.

\subsection{Tables that ran off the page}

Two layout faults, both of the kind that a compile without errors will happily
ship. The timing summary is 62 rows, and a plain \texttt{tabular} is a single
unbreakable box, so it ran off the bottom of the page and was clipped rather than
continued. It is now a \texttt{longtable}, which repeats its header and marks
continuations. The Nsight counter table has seven columns and ran off the right
margin instead; it is scaled to the text width. The appendix tables carry raw
metric names up to 49 characters with no break point, so they are set smaller
rather than scaled, because a \texttt{longtable} spanning pages cannot be put
inside a scaling box.

The general point is that LaTeX reports overfull boxes as warnings, not errors, so
a report can compile cleanly and still be unreadable. Checking the log for
overfull boxes is now part of judging whether a build is good, not just whether it
finished.

\end{document}
